\section{What the Fence Guarantees}
\label{sec:theory}

The durable state $\sigma$ contains scope epochs, grant flags and labelled copies.
Table~\ref{tab:lifecycle-states} lists the states of one submission. An authorization
log is an observation of intent, not a send. A \emph{handoff} $c$ is acceptance of a
positive byte prefix by the audited socket transport; a write may accept fewer bytes
than requested or none. The guarantee concerns each accepted prefix, including a
prefix of a request that never completes.

\begin{table}[t]\centering
\caption{Lifecycle of one submission, with the artifact's event names. The fence check shares the withdrawal lock (A3); earlier checks do not.}
\label{tab:lifecycle-states}\mtab
\begin{tabular}{@{}p{0.40\columnwidth}p{0.18\columnwidth}p{0.34\columnwidth}@{}}\toprule
\mhead State (event) & Check & Crash or loss here\\\midrule
Labelled copy stored (\texttt{COLLECT}) & none & label survives with the row (A2)\\
Job queued, label in row or extras (\texttt{ENQUEUE}) & enqueue check & job restored after reboot with its epochs\\
Batch prepared (\texttt{JOBSTART}, \texttt{AUTHORIZATION\_LOGGED}) & job-start check & nothing accepted; retry rereads policy\\
TLS record buffered (\texttt{TLS\_BUFFERED}) & none & nothing accepted\\
Prefix handed off (\texttt{WRITE\_ATTEMPT}, \texttt{BYTES\_ACCEPTED}) & fence, per write & count may be unlogged: unknown, never zero\\
Request received (receiver log) & none & accepted bytes may still arrive (Prop.~\ref{prop:recall})\\
Response received (\texttt{RESPONSE\_RECEIVED}) & none & lost response: retry may duplicate\\
\bottomrule\end{tabular}\end{table}

\begin{assumption}[Supported path]\label{as:sp}
(A1) Every supported byte is sent through Algorithm~\ref{alg:fence}, with one writer
per socket and no unmediated asynchronous sender below its check. For TLS, the checked bytes are application-record ciphertext: the request's conservative source union remains attached to every encrypted record and partial-record remainder. Encryption may precede checking; each subsequent socket write must recheck policy. Handshake traffic contains no request data. TLS early data, renegotiation and external writers are excluded.
(A2) Every persisted copy carries atomically stored, over-approximate source labels;
missing or invalid labels fail closed. Labels on retries retain the collection epochs.
(A3) Each bounded, nonblocking write and its policy check share a process/cross-process
lock with withdrawal and re-grant. A withdrawal commits its policy transaction before
acknowledgement. Policy state survives process failure; an unreadable store fails closed.
The model assumes working durable storage, not recovery from arbitrary power loss.
\end{assumption}

\begin{theorem}[Soundness of accepted prefixes]\label{thm:sound}
Under A1, A2 and A3, every source label of a supported prefix accepted at $c$ satisfies
$\neg\mathit{forbid}_{\sigma(c)}(s,e,p)$. Thus no prefix accepted after an acknowledged
withdrawal carries a source that remains invalid under the stated contract.
\end{theorem}
\begin{proof}
Consider one write and its preceding check in the same critical section. A3 excludes
policy changes between the check and every byte acceptance during that write. A2
includes every source in the checked label union, and the check rejects any forbidden
source. Hence every accepted prefix satisfies the rule. Crashes add no exception
(Table~\ref{tab:lifecycle-states}): a prefix accepted by an interrupted write is still
covered by its check, even if no count is logged, and after restart the next attempt
reacquires the lock and rereads durable policy, so it cannot reuse an earlier authorization.
\end{proof}

Exact labels avoid policy-induced false positives only when labels and policy are
available and valid. Over-approximation, fail-closed migration, timeouts and network
failure can still prevent permitted delivery. Soundness is not an availability claim.

\begin{theorem}[Race for earlier checks]\label{thm:race}
Suppose a check $q$ admits a source that a later withdrawal would invalidate for the
submission's purpose, and the path can still hand off that source at $c>q$ without
another effective check or cancellation. Then a withdrawal in $(q,c)$ yields a forbidden
handoff. If a single such withdrawal $W$ is independent of fixed handoff offsets
$c_1<\cdots<c_k$, all carrying sources invalidated by $W$, with no intervening re-grant and $n_i>0$ newly handed-off source items,
\begin{equation}
\begin{aligned}
\Pr[\text{violation}]&=\Pr[q<W<c_k],\\
\mathbb E[\#\text{forbidden item handoffs}]&=\sum_i n_i\Pr[q<W<c_i].
\end{aligned}\label{eq:race}
\end{equation}
These equations concern that schedule model; cancellations or retries that depend on
$W$ require a different model.
\end{theorem}
\begin{proof}
Place the applicable withdrawal between the passing check and the unprotected write.
The path still sends, by hypothesis. Summing the corresponding indicator variables
gives the expectation; the last interval contains the earlier intervals.
\end{proof}

\begin{proposition}[Information needed by the check]\label{prop:nec}
A declared-scope check can admit an invalid source when a withdrawal invalidates the
actual source scope while the declared scope still passes. A current boolean alone
cannot distinguish old and new data after withdrawal and re-grant when the contract
invalidates pending old data. Source lineage and retained epochs supply these distinctions;
equivalent representations are possible.
\end{proposition}

\begin{proposition}[No recall]\label{prop:recall}
A prefix accepted before withdrawal may arrive after acknowledgement. A device-side
fence cannot guarantee recall of bytes already accepted by the transport.
\end{proposition}

Each retry is checked afresh, but the mechanism provides neither server-side deletion nor
exactly-once delivery. The stream-wrapper evaluation's pre-write sequence markers cannot establish which
prefixes were outstanding; the channel wrapper records acceptance and receipt separately.
